\documentclass[journal,twoside]{IEEEtran}
\usepackage{cite}
\usepackage{amsmath,amssymb,amsfonts}
\usepackage{algorithm}
\usepackage{algpseudocode}
\usepackage{booktabs}
\usepackage{graphicx}
\usepackage{textcomp}
\usepackage{multirow}
\usepackage{array}
\usepackage{xcolor}
\usepackage{url}

% Visible placeholder for every number that depends on the new (MPCE) runs.
% (an empty argument, used in table cells, prints just [TBD])
\newcommand{\TBD}[1]{\textcolor{red}{[TBD\ifx\relax#1\relax\else: #1\fi]}}
% Figure from the analysis pipeline; a framed placeholder is printed while the file does not exist yet.
% #1 width, #2 file (with extension), #3 placeholder height
\newcommand{\pipefig}[3]{\IfFileExists{#2}{\includegraphics[width=#1]{#2}}%
{\fbox{\parbox[c][#3][c]{\dimexpr#1-2\fboxsep-2\fboxrule\relax}{\centering\TBD{figure {\ttfamily\detokenize{#2}} not yet generated}}}}}

\markboth{Journal of Modern Power Systems and Clean Energy, VOL. XX, NO. XX, XXXX}
{SOLANKI \MakeLowercase{\textit{et al.}}: TWO-PHASE PARTICLE SWARM AND VARIABLE NEIGHBORHOOD SEARCH FOR WIND FARM LAYOUT OPTIMIZATION}

\begin{document}

\title{A Two-Phase Particle Swarm and Variable Neighborhood Search Algorithm for Continuous Wind Farm Layout Optimization}
\author{Prince~Solanki, Pulkit~Dwivedi, Vanita~Garg, and Vinay~Shukla
\thanks{Manuscript received: \TBD{date}; revised: \TBD{date}; accepted: \TBD{date}.}
\thanks{\TBD{funding statement, e.g., ``This work did not receive any specific grant from funding agencies.''}}
\thanks{P. Solanki \TBD{mark corresponding author} is with the Department of Mathematics, School of Sciences, IILM University, Greater Noida, Uttar Pradesh, India; he was previously with the Department of Mathematics, Indian Institute of Technology Roorkee, Roorkee, India (e-mail: prince@ma.iitr.ac.in; prince@iilm.edu).}
\thanks{P. Dwivedi is with the Department of Computer Science and Engineering, Bennett University, Greater Noida, Uttar Pradesh, India (e-mail: pdwivedi19900@gmail.com; pulkit.dwivedi@bennett.edu.in).}
\thanks{V. Garg is with Amity University, Noida, Uttar Pradesh, India (e-mail: vanitagarg16@gmail.com).}
\thanks{V. Shukla is with the School of Mathematical Sciences, Shanghai Jiao Tong University, Shanghai 200240, China, and also with the Department of Mathematics, School of Computer Science Engineering and Technology, Bennett University, Greater Noida, India (e-mail: vinayshukla4321@gmail.com; vnshukla01@sjtu.edu.cn).}}

\maketitle

\begin{abstract}
The wind farm layout optimization problem (WFLOP) asks where a prescribed number of turbines should be placed to minimize wake losses while every turbine stays inside the site and keeps a minimum distance from the others. This paper proposes PSO-VNS, a two-phase algorithm for the continuous WFLOP. Particle swarm optimization (PSO) with the constriction coefficients of Clerc and Kennedy explores the layout space for the first half of the evaluation budget, and the basic variable neighborhood search (VNS) then refines the global best layout of the swarm. Both phases minimize one penalized objective based on the Jensen wake model and direction-dependent Weibull wind data. PSO-VNS is compared with PSO, a salp swarm variant of the same design (SSA-VNS), the salp swarm algorithm (SSA), the Laplacian SSA (LX-SSA), differential evolution, VNS and a multistart SLSQP solver on 68 benchmark cases, with 30 seed-paired runs of 6,030 objective evaluations per method and case. PSO-VNS obtains the best average rank (\TBD{avg.\ rank}) and \TBD{summary of Holm--Wilcoxon outcome and feasibility rate}. The ablation shows that \TBD{contribution of the VNS phase and of the swarm start, incl.\ vs.\ random multistart}, that SSA and LX-SSA are weaker first phases than a convergent PSO, and that the Laplace step of LX-SSA gives no gain. We also find that the widely used PSO setting $w=0.7$, $c_1=c_2=2$, which lies outside the convergence region of PSO, understates PSO as a baseline. Tests with feasibility-preserving initialization, budgets of up to 120,030 evaluations, a Horns Rev~1 block and IEA Wind Task~37 Case Study~1 show \TBD{one-sentence summary}.
\end{abstract}

\begin{IEEEkeywords}
Wind farm layout optimization, particle swarm optimization, constriction coefficient, variable neighborhood search, hybrid metaheuristic, wake effect.
\end{IEEEkeywords}

\section*{Nomenclature}
\addcontentsline{toc}{section}{Nomenclature}
\begin{IEEEdescription}[\IEEEusemathlabelsep\IEEEsetlabelwidth{$\xi_1,\xi_2,\xi_3$}]
\item[$B$] Budget (number of objective calls).
\item[$c_1$, $c_2$] Cognitive and social acceleration coefficients of PSO.
\item[$C_T$] Thrust coefficient.
\item[$D$, $R$] Rotor diameter and rotor radius ($D=2R$).
\item[$d_{ij}$] Downstream distance of turbine $i$ behind turbine $j$.
\item[$\Delta_{ij}$, $\Delta_i$] Velocity deficit at turbine $i$ due to turbine $j$; combined deficit at turbine $i$.
\item[$E(P_i)$] Expected power of turbine $i$ (benchmark units).
\item[$F_p$] Penalized wake loss (minimized).
\item[$f(s)$] Power curve.
\item[$\mathbf G$, $\mathbf P^i$] Global best and personal best of particle $i$ (PSO).
\item[$\mathbf H$] Food source (best layout of the salp swarm).
\item[$K$, $K_{\rm G}$] Jensen wake-spreading constant; growth rate of the Gaussian wake.
\item[$k$, $k_{\max}$] Index and number of VNS neighborhoods.
\item[$\ell_{ij}$, $\ell_{\min}$] Distance between turbines $i$ and $j$; minimum spacing.
\item[$N$, $N_p$] Number of turbines; population size.
\item[$N_\theta$, $N_s$] Numbers of wind-direction and wind-speed bins.
\item[$p_l$] Probability of the $l$th wind-direction bin.
\item[$r$] Farm radius.
\item[$s$] Wind speed.
\item[$t$, $T$, $T_1$] Iteration counter, number of swarm iterations, and number of swarm iterations in Phase~1.
\item[$\mathbf V^i$] Velocity of particle $i$.
\item[$w$] Inertia weight of PSO.
\item[$\mathbf X$] Layout vector $(x_1,y_1,\ldots,x_N,y_N)\in\mathbb R^{2N}$.
\item[$x_i$, $y_i$] Coordinates of turbine $i$.
\item[$\alpha$, $\beta_{ij}$] Half-cone angle of the wake; wake angle of turbine $i$ with respect to turbine $j$.
\item[$\delta_k$] Outer radius of the $k$th VNS neighborhood.
\item[$\zeta$, $\psi$] Weibull shape and scale parameters.
\item[$\theta$] Wind direction (direction the wind blows towards, counter-clockwise from east).
\item[$\lambda_1$, $\lambda_0$] Slope and intercept of the linear power curve.
\item[$\mu$] Penalty factor.
\item[$\boldsymbol\rho_1$, $\boldsymbol\rho_2$] Vectors of uniform random numbers in the PSO update.
\item[$\phi$, $\chi$] Location and scale of the Laplace step (LX-SSA).
\item[$\xi_1,\xi_2,\xi_3$] SSA leader coefficient and uniform random numbers.
\item[$\omega$] Fraction of the budget given to Phase~1.
\end{IEEEdescription}

\section{Introduction}

Wind energy has become one of the main sources of low-carbon electricity, and the placement of the turbines, fixed at the planning stage, has a lasting effect on the energy a farm delivers. The wind farm layout optimization problem (WFLOP) addresses this decision. An upstream turbine leaves behind a wake of slower and more turbulent air that reduces the power of the turbines downstream; the wake recovers with distance, but the site is limited and the turbines must keep a minimum distance from each other. The objective is nonlinear, the wakes couple the positions of all turbines, and the feasible region is bounded by many nonconvex spacing constraints, which has made WFLOP a popular test bed for metaheuristics.

Early studies placed turbines on a grid with genetic algorithms (GAs)~\cite{Mosetti1994,Grady2005}. Kusiak and Song~\cite{Kusiak2010} introduced the continuous, wind-distribution-based model that we also use, and it has since been solved with ant colony optimization and particle filtering~\cite{Eroglu2012,Eroglu2013} and with biogeography-based optimization~\cite{Bansal2017}. GA, particle swarm optimization (PSO), differential evolution (DE) and simulated annealing were compared for offshore layouts in~\cite{Elkinton2008}, PSO has been applied to reduce wake losses~\cite{Asaah2021}, and related design problems such as the collector system and the electrical layout of offshore farms have been reviewed and optimized in this journal~\cite{Srikakulapu2018,Hou2019}. Later work combined GAs with problem-specific information~\cite{Ju2019} or Monte Carlo tree search~\cite{Bai2022}, refined GA layouts with deterministic local methods~\cite{Nagpal2021}, developed a variable neighborhood search (VNS) heuristic for large offshore farms~\cite{Cazzaro2022}, and introduced pseudo-gradients~\cite{Quaeghebeur2021}, Gaussian wake models~\cite{Tao2020} and data-driven or CFD-based surrogates~\cite{Yang2023,Wang2024,Li2025}. The IEA Wind Task~37 case study~\cite{Baker2019} showed that, even with a common wake model, different optimizers reach clearly different layouts.

Many of these studies propose a new or modified metaheuristic and compare it with reference methods, among which PSO is one of the most common. Such comparisons are only as informative as the configuration of the references. PSO converges only if its inertia weight $w$ and acceleration coefficients $c_1$, $c_2$ lie in a known region~\cite{Clerc2002,Poli2009}, and the setting $w=0.7$, $c_1=c_2=2$, which is frequently used for PSO baselines and was used in the earlier version of this study, lies outside it. With the constriction coefficients of Clerc and Kennedy~\cite{Clerc2002}, plain PSO outperforms the salp swarm algorithm (SSA)~\cite{Mirjalili2017}, the Laplacian SSA (LX-SSA)~\cite{Solanki2023} proposed earlier by two of the authors, and an SSA-VNS hybrid that we first developed (Section~\ref{sec:results}). A baseline outside its convergence region can therefore make a new method look stronger than it is; we report this as a finding about the configuration of baselines, not about any particular study.

Swarms and trajectory methods complement each other. Once a swarm has contracted around its best layout, the remaining evaluations bring small gains, because the final refinement requires coordinated moves of single turbines along the spacing constraints, whereas VNS~\cite{Mladenovic1997,Hansen2001} refines one layout efficiently but depends on where it starts. A swarm-then-VNS design helps only if the swarm provides a better start than cheaper alternatives, and SSA did not: SSA-VNS was statistically tied with a control that spends the same evaluations on random layouts (Section~\ref{sec:ablation}). A convergent PSO, which contracts around a good region within its share of the budget, is the natural alternative.

We therefore propose PSO-VNS: PSO with constriction coefficients explores the layout space for the first half of the budget, and the basic VNS refines the global best layout of the swarm. Both components are standard; the contribution lies in their combination for the constrained WFLOP and in a controlled evaluation of which ingredients matter. We keep the Jensen wake model, Weibull wind data and the Kusiak--Song power model for comparability with the literature and examine these choices separately. Our contributions are as follows:
\begin{enumerate}
\item We propose PSO-VNS, a two-phase algorithm for the continuous WFLOP with explicit boundary and minimum-spacing constraints, in which a convergent PSO and the basic VNS share one penalized objective (Section~\ref{sec:hybrid}).
\item We compare PSO-VNS with PSO, SSA-VNS, SSA, LX-SSA, DE, VNS and a gradient-based multistart SLSQP solver on 68 benchmark cases at equal numbers of objective evaluations with 30 seed-paired runs, and isolate in an ablation the VNS phase, the value of the swarm start (against VNS alone and a random-multistart control), the choice of the first-phase swarm and the Laplace step of LX-SSA (Sections~\ref{sec:results} and~\ref{sec:ablation}).
\item We show that the PSO setting $w=0.7$, $c_1=c_2=2$ violates the convergence condition of PSO and quantify how much it understates PSO (Sections~\ref{sec:pso} and~\ref{sec:psosetting}).
\item We test the methods with feasibility-preserving initialization, budgets of up to 120,030 evaluations, a 16-turbine block of the Horns Rev~1 offshore farm, the 16- and 36-turbine scenarios of IEA Wind Task~37 Case Study~1 (IEA37), and alternative power-curve, wake, spacing and boundary assumptions (Sections~\ref{sec:beyond} and~\ref{sec:robustness}).
\end{enumerate}

Section~\ref{sec:model} describes the model, Section~\ref{sec:prelim} the swarm methods, and Section~\ref{sec:hybrid} PSO-VNS. Sections~\ref{sec:setup}--\ref{sec:robustness} give the setup and results, Section~\ref{sec:limits} the limitations, and Section~\ref{sec:conclusion} concludes the paper.

\section{Problem Modeling}
\label{sec:model}

We use the following definitions and assumptions.
\begin{enumerate}
    \item \textbf{Turbines and layout:} The number of turbines $N$ is prescribed in every run, and only the $2N$ coordinates $(x_i,y_i)$, measured from the center of the farm, are optimized; results for several $N$ form a sweep, and $N$ is not optimized. All turbines are identical.
    \item \textbf{Wind:} For a wind direction $\theta$, the wind speed $s$ follows a Weibull distribution with density
    $p_s(s;\zeta,\psi)=(\zeta/\psi)(s/\psi)^{\zeta-1}\exp[-(s/\psi)^\zeta]$, whose shape $\zeta=\zeta(\theta)$ and scale $\psi=\psi(\theta)$ depend on the direction. Throughout the paper, $\theta$ is the direction \emph{towards} which the wind blows, measured counter-clockwise from east, so that turbine $i$ lies downstream of turbine $j$ when $(x_i-x_j)\cos\theta+(y_i-y_j)\sin\theta>0$. Meteorological directions $\varphi$ (the direction the wind comes \emph{from}, clockwise from north), as given for Horns Rev~1 and IEA37, are converted by $\theta=270^\circ-\varphi$.
    \item \textbf{Constraints:} Any two turbines must be at least $\ell_{\min}=4D=8R$ apart, $\ell_{ij}=\sqrt{(x_i - x_j)^2 + (y_i - y_j)^2} \geq \ell_{\min}$, the value of the benchmark studies~\cite{Kusiak2010,Bansal2017} (Section~\ref{sec:spacing} studies $5D$ and $6D$). The farm is circular~\cite{Kusiak2010,Bansal2017}, $x_i^2+y_i^2\le r^2$, with $r=500$, 750 and 1000~m, and the turbines are not restricted to a grid.
    \item \textbf{Objective:} We use the Jensen wake model~\cite{Jensen1983,Katic1986} and the power model of~\cite{Kusiak2010}, and maximize the expected power of the farm, i.e., minimize the wake losses, subject to both constraints.
\end{enumerate}

\subsection{Wake Effect Model}
In the Jensen model~\cite{Jensen1983,Katic1986}, the wake behind a rotor expands linearly with the downstream distance. With the wake-spreading constant $K$, the thrust coefficient $C_T$ and the distance $d_{ij}$ of turbine $i$ behind turbine $j$, the velocity deficit is
\begin{equation}
    \Delta_{ij} = 1-\frac{s_{\text{down}}}{s_{\text{up}}}=\frac{1 - \sqrt{1 - C_T}}{\left(1 + K d_{ij}/R\right)^2}.
    \label{eq:deficit}
\end{equation}
The wake of turbine $j$ is a half cone with half-angle $\alpha=\arctan K$ whose virtual vertex lies $R/K$ upstream of the rotor. Turbine $i$ lies in it if the wake angle
$\beta_{ij}=\cos^{-1}\!\bigl(u_{ij}/\sqrt{u_{ij}^2+v_{ij}^2}\bigr)$ satisfies $\beta_{ij}<\alpha$, where
\begin{align*}
u_{ij}&=(x_i - x_j)\cos\theta + (y_i - y_j)\sin\theta + R/K,\\
v_{ij}&=-(x_i - x_j)\sin\theta + (y_i - y_j)\cos\theta ,
\end{align*}
$v_{ij}$ is the lateral offset from the wake axis and $d_{ij}=\lvert u_{ij}-R/K\rvert$. The deficits are combined by the root-sum-square rule,
\begin{equation}\label{eq5}
\Delta_i(\theta) = \Bigl( \textstyle\sum_{j \ne i,\ \beta_{ij} < \alpha} \Delta_{ij}^2 \Bigr)^{1/2} .
\end{equation}
Within a direction bin, $\Delta_i(\theta)$ reduces the free-stream speed by a constant factor. Since $aS\sim\mathrm{Weibull}(\zeta,a\psi)$ for $a>0$ when $S\sim\mathrm{Weibull}(\zeta,\psi)$, the local wind speed keeps the shape parameter and has the scale parameter
\begin{equation}
\psi_i(\theta) = \psi(\theta)\left[1-\Delta_i(\theta)\right],\quad i=1,\ldots,N. \label{eq:weibull-scale}
\end{equation}
Schematics of the wake model and the wind-direction convention are given in Figs.~S1--S3 of the supplementary material.

\subsection{Power Model}
\label{sec:powermodel}
We use the power curve of the Kusiak--Song benchmark~\cite{Kusiak2010}: $f(s)=0$ for $s<s_{\text{cut-in}}$, $f(s)=\lambda_1 s+\lambda_0$ for $s_{\text{cut-in}}\le s\le s_{\text{rated}}$ and $f(s)=P_{\text{rated}}$ above, with $s_{\text{cut-in}}=3.5$~m/s, $s_{\text{rated}}=14$~m/s, $P_{\text{rated}}=1500$~kW, $\lambda_1=140.86$ and $\lambda_0=-500$. We use this curve exactly as published, although its linear part is slightly negative just above the cut-in speed and reaches 1472~kW at the rated speed, where it jumps to 1500~kW; the benchmark has no cut-out speed. The expected power of turbine $i$,
\begin{equation}\label{eq8}
E(P_i) = \int_0^{360} p_\theta(\theta) \int_0^{\infty} f(s)\,
       p_s\bigl(s; \zeta(\theta), \psi_i(\theta)\bigr) \, ds \, d\theta ,
\end{equation}
is evaluated, as in~\cite{Kusiak2010}, by a Riemann sum over $N_\theta$ direction bins with probabilities $p_l$ and $N_s$ speed bins between $s_{\text{cut-in}}$ and $s_{\text{rated}}$ (the discrete form is given in the supplementary material). Like the earlier studies, we keep the bin width in degrees as a factor of the direction sum. Because all 24 direction bins are $15^\circ$ wide and $\sum_l p_l=1$, the objective equals 15 times the expected farm power in kW, and $\mathrm{AEP}\,[\mathrm{MWh}] = (\text{objective}/15)\times 8.76$; e.g., the wake-free objective of one turbine under Wind Data Set~I, $14045.74$, corresponds to 936.38~kW or 8202.7~MWh per year.

\subsection{Constraint Handling}
\label{sec:constraints}
After every position update, each coordinate is clipped to the square that encloses the farm, $x_i\leftarrow\min\{\max\{x_i,-r\},r\}$ and likewise for $y_i$ (box repair). The circular boundary and the spacing are handled by a penalty. With the violations
$g^{\rm s}_{ij}=\ell_{\min}-\ell_{ij}$ and $g^{\rm b}_{i}=x_i^2+y_i^2-r^2$,
a layout is feasible if $g^{\rm s}_{ij}\le0$ and $g^{\rm b}_i\le0$ for all $i$ and $j$. All penalty-based optimizers minimize the penalized wake loss
\begin{equation}
\begin{split}
F_p(\mathbf X)={}&\Bigl(N E_{\rm ideal}-\sum_{i=1}^{N}E(P_i)\Bigr)+\sum_{i:\,g^{\rm b}_i>0}\bigl(1+\mu g^{\rm b}_i\bigr)^2\\
&+\sum_{i<j:\,g^{\rm s}_{ij}>0}\bigl(1+\mu g^{\rm s}_{ij}\bigr)^2,\qquad \mu=10^{10},
\end{split}
\label{eq:penalty}
\end{equation}
where $\mathbf X=(x_1,y_1,\ldots,x_N,y_N)$ and $N E_{\rm ideal}$ is the wake-free objective, so that the first term is the wake loss reported in our tables. Any violation larger than about $10^{-7}$ (m or m$^2$) makes the penalty larger than the largest possible wake loss. Comparing $F_p$ values therefore gives a feasibility-first selection: a feasible layout is always preferred to an infeasible one, two feasible layouts are compared by their wake loss, and two infeasible layouts by their total penalty.

\section{Preliminaries: Swarm Methods for Phase~1}
\label{sec:prelim}
All swarm methods below keep a population of $N_p$ candidate vectors $\mathbf X^i\in\mathbb R^{n}$ with components $X^i_m$, $m=1,\ldots,n$, bounds $lb_m$, $ub_m$, and cost $N_p$ objective calls per iteration unless stated otherwise.

\subsection{Particle Swarm Optimization with Constriction Coefficients}
\label{sec:pso}
In PSO~\cite{Kennedy1995}, particle $i$ has a position $\mathbf X^i$, a velocity $\mathbf V^i$ and a personal best $\mathbf P^i$, and the swarm keeps the global best $\mathbf G$. In the inertia-weight form~\cite{Shi1998}, a particle moves by
\begin{align}
\mathbf V^i&\leftarrow w\,\mathbf V^i+c_1\boldsymbol\rho_1\circ\bigl(\mathbf P^i-\mathbf X^i\bigr)+c_2\boldsymbol\rho_2\circ\bigl(\mathbf G-\mathbf X^i\bigr),\label{eq:pso-v}\\
\mathbf X^i&\leftarrow\mathbf X^i+\mathbf V^i,\label{eq:pso-x}
\end{align}
where $\boldsymbol\rho_1,\boldsymbol\rho_2$ are vectors of independent $U(0,1)$ numbers, drawn anew for every move, and $\circ$ is the componentwise product. Clerc and Kennedy~\cite{Clerc2002} multiply the bracket $\mathbf V^i+\varphi_1\boldsymbol\rho_1\circ(\mathbf P^i-\mathbf X^i)+\varphi_2\boldsymbol\rho_2\circ(\mathbf G-\mathbf X^i)$ by the constriction factor $\kappa=2/\lvert2-\varphi-\sqrt{\varphi^2-4\varphi}\rvert$, $\varphi=\varphi_1+\varphi_2$. The standard choice $\varphi_1=\varphi_2=2.05$ gives $\kappa=0.7298$, which is \eqref{eq:pso-v} with
\begin{equation}
w=0.7298,\qquad c_1=c_2=\kappa\varphi_1=1.49618 .
\label{eq:constriction}
\end{equation}

When the personal and global bests stay fixed (stagnation), the mean and the variance of the particle positions converge if $\lvert w\rvert<1$ and~\cite{Poli2009}
\begin{equation}
c_1+c_2<\frac{24\,(1-w^2)}{7-5w}.
\label{eq:poli}
\end{equation}
For the constriction setting~\eqref{eq:constriction}, the right-hand side equals 3.35 and $c_1+c_2=2.99$ satisfies~\eqref{eq:poli}. For $w=0.7$ the bound is 3.50, so the common setting $c_1=c_2=2$ ($c_1+c_2=4$) violates it: the variance of the positions then grows even when the bests no longer change, and the swarm keeps sampling far from its best layouts instead of refining them. In our implementation, the velocities start at zero and are not bounded, the positions are box-repaired after each move, and $\mathbf P^i$ and $\mathbf G$ are updated immediately after each particle is evaluated.

\subsection{Salp Swarm Algorithm and Its Laplacian Variant}
\label{sec:ssa}
The SSA~\cite{Mirjalili2017} sorts the population by the objective, keeps the best solution found so far as the food source $\mathbf H$, and splits the population into a leader half and a follower half. A leader $i\le N_p/2$ updates each coordinate around the food source,
\begin{equation}\label{eq:leader}
X_m^i = H_m\pm\xi_1\big((ub_m-lb_m)\xi_2+lb_m\big),
\end{equation}
with the sign chosen with probability $1/2$ by $\xi_3$, where $\xi_2,\xi_3\sim U(0,1)$ and $\xi_1=2\exp[-(4t/T)^2]$ decreases with the iteration $t$ out of $T$. A follower $i>N_p/2$ moves to the midpoint between itself and the preceding salp, $X_m^i\leftarrow\tfrac{1}{2}(X_m^i+X_m^{i-1})$. LX-SSA~\cite{Solanki2023} gives each follower a second, Laplace-based candidate~\cite{Deep2007}
\begin{align}
Y_m^i&=X_m^i+\gamma\left(H_m-X_m^i\right),\nonumber\\
\gamma&=\phi-\chi\,\mathrm{sgn}\bigl(z-\tfrac12\bigr)\ln\bigl(1-2\lvert z-\tfrac12\rvert\bigr),
\label{eq:laplace}
\end{align}
with $z\sim U(0,1)$, i.e., a Laplace variable with location $\phi=0$ and scale $\chi=1$, and the follower keeps the better of its two candidates. One LX-SSA iteration therefore costs $2N_p$ calls. The comparison of SSA and LX-SSA isolates the Laplace step. The pseudocode of both methods is given in the supplementary material.

\section{Proposed PSO-VNS Algorithm}
\label{sec:hybrid}

PSO-VNS combines a swarm that explores the layout space and settles around a good region with a trajectory method that refines one layout. The first half of the budget is spent by a convergent PSO, and the basic VNS refines the global best layout of the swarm in the second half.

\subsection{Application to the WFLOP}
For a prescribed turbine count $N$, a candidate is the layout vector $\mathbf X\in\mathbb R^{n}$, $n=2N$, with bounds $[-r,r]^{n}$. Every candidate is box-repaired, its deficits $\Delta_i(\theta)$ and scale parameters $\psi_i(\theta)$ are computed for every direction bin by \eqref{eq5} and \eqref{eq:weibull-scale}, and $F_p$ is evaluated by \eqref{eq:penalty}. Both phases use $F_p$ in every comparison, best update and acceptance step, so the feasibility-first rules of Section~\ref{sec:constraints} hold throughout; in particular, $\mathbf P^i$ and $\mathbf G$ move towards feasible layouts as soon as the swarm finds them.

\subsection{Phase 1: Exploration by Constriction PSO}
For a budget of $B$ objective calls and a split fraction $\omega\in(0,1)$, PSO with the coefficients \eqref{eq:constriction} runs with population $N_p$ for
\[
T_1=\operatorname{round}\!\left(\frac{\omega B-N_p}{N_p}\right)
\]
iterations, since the initial swarm costs $N_p$ calls and every iteration another $N_p$. Because the constriction PSO has no iteration-dependent parameter, Phase~1 of PSO-VNS coincides, for a given seed, with the first $N_p(T_1+1)$ calls of a PSO run with the same budget; PSO-VNS and PSO differ only in how they spend the rest of the budget. The global best $\mathbf G$ at the end of Phase~1 is the starting point of Phase~2.

\subsection{Phase 2: Refinement by Basic VNS}
Phase~2 is the basic VNS of~\cite{Mladenovic1997,Hansen2001} in the continuous form of~\cite{Mladenovic2008}, applied to $F_p$. It uses $k_{\max}=5$ nested $\ell_\infty$ neighborhoods of the whole layout,
\[
\mathcal N_k(\mathbf X)=\{\mathbf X':\delta_{k-1}<\lVert \mathbf X'-\mathbf X\rVert_\infty\le\delta_k\},\qquad \delta_k=0.1\,k\,r .
\]
In the shaking step, a point is drawn uniformly from $\mathcal N_k(\mathbf X)$, so that all $2N$ coordinates move, and box repair is applied. The shaken layout is then improved by a complete best-improvement local search. Since the top-hat Jensen wake makes $F_p$ discontinuous, we use a derivative-free compass search: it evaluates all $2n=4N$ moves $\pm h\,\mathbf e_m$ along the coordinate axes, moves to the best improving one, and halves the step $h$ (starting from $0.05r$) when no move improves, until $h<10^{-3}r$. If the resulting local minimum is better than the incumbent, it is accepted and $k$ is reset to 1; otherwise $k$ is increased by one and returns to 1 after $k_{\max}$. Phase~2 first applies the local search to $\mathbf G$ and then repeats shaking, local search and neighborhood change until the budget $B$ is used up (Algorithm~\ref{alg:hybrid}).

\begin{algorithm}[!t]
\caption{PSO-VNS for a prescribed turbine count $N$}
\label{alg:hybrid}
\begin{algorithmic}[1]
\Require budget $B$, split $\omega$, population $N_p$, $w$, $c_1$, $c_2$, farm radius $r$, $k_{\max}$
\State $T_1\gets\operatorname{round}((\omega B-N_p)/N_p)$ \Comment{Phase 1}
\State Draw and evaluate $N_p$ particles $\mathbf X^i$; $\mathbf V^i\gets\mathbf 0$; $\mathbf P^i\gets\mathbf X^i$; $\mathbf G\gets\arg\min_i F_p(\mathbf P^i)$
\For{$t=1$ to $T_1$}
    \For{particle $i=1,\ldots,N_p$}
        \State Update $\mathbf V^i$, $\mathbf X^i$ by \eqref{eq:pso-v}, \eqref{eq:pso-x}; box repair
        \State Evaluate $F_p(\mathbf X^i)$; update $\mathbf P^i$ and $\mathbf G$
    \EndFor
\EndFor
\State $\mathbf X\gets\textsc{LocalSearch}(\mathbf G)$; $k\gets1$ \Comment{Phase 2}
\While{fewer than $B$ objective calls have been used}
    \State $\mathbf X'\gets$ uniform random point of $\mathcal N_k(\mathbf X)$; box repair
    \State $\mathbf X'\gets\textsc{LocalSearch}(\mathbf X')$
    \If{$F_p(\mathbf X')<F_p(\mathbf X)$}
        \State $\mathbf X\gets \mathbf X'$; $k\gets1$
    \Else
        \State $k\gets k\bmod k_{\max}+1$
    \EndIf
\EndWhile
\State \Return best layout evaluated
\Statex
\Function{LocalSearch}{$\mathbf X'$} \Comment{compass search}
    \State $h\gets0.05r$
    \While{$h\ge10^{-3}r$}
        \State $\mathbf X^\star\gets\arg\min F_p$ over the $2n$ points $\mathbf X'\pm h\,\mathbf e_m$ (box-repaired)
        \If{$F_p(\mathbf X^\star)<F_p(\mathbf X')$} \State $\mathbf X'\gets \mathbf X^\star$ \Else \State $h\gets h/2$ \EndIf
    \EndWhile
    \State \Return $\mathbf X'$
\EndFunction
\end{algorithmic}
\end{algorithm}

\subsection{Parameters and Cost}
PSO-VNS uses $N_p=30$, the coefficients \eqref{eq:constriction} and the VNS settings above, and it adds a single parameter, the split $\omega$. With the default $\omega=0.5$ and $B=6{,}030$, PSO runs for $T_1=100$ iterations and receives $30+100\times30=3{,}030$ calls, and VNS the remaining 3,000. None of these values was tuned for PSO-VNS; a study of the split (25\%, 50\% and 75\%) was carried out only for SSA-VNS and is reported in the supplementary material. The cost of both phases is dominated by the objective evaluations, so for a budget of $B$ calls PSO-VNS has the same $\mathcal O(B\,N^2N_\theta N_s)$ cost as the other methods; the PSO update, box repair, shaking and each compass move are $\mathcal O(N)$ and the spacing check $\mathcal O(N^2)$. Because the random-number generator is seeded once and Phase~1 draws the initial population first, PSO-VNS starts from the same initial population as all other methods for a given seed, which keeps the comparisons paired.

\section{Experimental Setup}
\label{sec:setup}

\subsection{Benchmark Cases}
We use the two wind data sets of the Kusiak--Song benchmark~\cite{Kusiak2010,Bansal2017}, listed in the supplementary material. In Wind Data Set~I, the Weibull parameters are the same in all directions ($\zeta=2$, $\psi=13$) and the wind is strongly directional (probabilities 0.2 and 0.6 for the bins $[75^\circ,90^\circ)$ and $[90^\circ,105^\circ)$, 0.01 for most others); Wind Data Set~II keeps $\zeta=2$, but $\psi$ varies from 2.6 to 10~m/s with the direction and the wind is spread over more directions. We use $N_\theta=24$ direction bins of $15^\circ$ and $N_s=21$ speed intervals of 0.5~m/s. Each data set is combined with three circular farms of radius $r=500$, 750 and 1000~m, with $N=2,\ldots,10$, $2,\ldots,12$ and $2,\ldots,15$ turbines, respectively, which gives $2\times(9+11+14)=68$ cases. The turbines have $R=38.5$~m, $\ell_{\min}=4D=308$~m, $C_T=0.8$ and $K=0.075$.

\subsection{Methods Compared}
The main comparison includes PSO-VNS and seven reference methods, all evaluated with the objective~\eqref{eq:penalty}, the bounds $[-r,r]^{2N}$, population 30 where applicable, and the same budget:
\begin{itemize}
\item \emph{PSO-VNS} (proposed): Algorithm~\ref{alg:hybrid} with $\omega=0.5$, i.e., 100 PSO iterations (3,030 calls) followed by 3,000 calls of VNS.
\item \emph{PSO}~\cite{Kennedy1995,Clerc2002}: the coefficients \eqref{eq:constriction} and 200 iterations. The setting $w=0.7$, $c_1=c_2=2$ of the earlier version of this study is discussed in Section~\ref{sec:psosetting}.
\item \emph{SSA-VNS}: Algorithm~\ref{alg:hybrid} with SSA in Phase~1 (100 SSA iterations, 3,030 calls, $\xi_1$ decreasing over these 100 iterations).
\item \emph{SSA}~\cite{Mirjalili2017}: 200 iterations; \emph{LX-SSA}~\cite{Solanki2023}: $\phi=0$, $\chi=1$ and 100 iterations.
\item \emph{DE}~\cite{Storn1997}: DE/rand/1/bin with scale factor 0.5, crossover rate 0.9 and 200 iterations.
\item \emph{VNS}: Phase~2 of Algorithm~\ref{alg:hybrid}, started from the best member of the initial population.
\item \emph{Multistart SLSQP (MS-SLSQP)}: the SLSQP solver of SciPy~\cite{Kraft1988} minimizes the wake loss with the spacing and boundary constraints imposed explicitly and forward-difference gradients (1-m step), starting from the members of the initial population in the order of their $F_p$ values and restarting until the budget is used up.
\end{itemize}
The ablation adds two variants. \emph{RS-VNS} is a random-multistart control: Phase~1 is replaced by pure random sampling with the same number of calls ($\operatorname{round}(\omega B)=3{,}015$ layouts, including the initial population), and the best sampled layout is refined by the same VNS. \emph{LX-SSA-VNS} uses LX-SSA in Phase~1 ($T_1=\operatorname{round}((\omega B-N_p)/(2N_p))=50$ iterations, 3,030 calls). SSA, LX-SSA, PSO and DE use the authors' original implementation, with the coefficients \eqref{eq:constriction} for PSO; VNS, RS-VNS, MS-SLSQP and the three hybrids were implemented for this study, and the hybrids call the original swarm code in Phase~1.

Every run of the main comparison has a budget of exactly 6,030 objective calls, counted by a wrapper around the objective function (including the initial population and, for MS-SLSQP, the finite-difference gradient calls), i.e., 200 iterations of PSO, SSA and DE and 100 of LX-SSA. Each method is run with the seeds 1--30; every optimizer seeds the random-number generator and draws its initial population first, so runs with the same seed start from the \emph{same} initial population for all methods and are paired.

\subsection{Tests Beyond the Benchmark}
\label{sec:setup-beyond}
\paragraph*{Initialization and budget} In the main comparison, the initial layouts are drawn uniformly in the bounding square and are usually infeasible. As an alternative, each initial layout is drawn inside the site and pushed apart by minimizing a smooth packing penalty with L-BFGS-B until the constraints hold, without objective calls. All ten methods (the eight of the main comparison, RS-VNS and LX-SSA-VNS) are run with this initialization at 6,030 calls, and with random initialization at 6,030, 30,030 and 120,030 calls, on the six largest cases (the largest $N$ of each farm and data set) and the Horns Rev~1 block, with 30 seeds (10 for Horns Rev~1 at 120,030 calls). All iteration counts scale with the budget; e.g., PSO-VNS uses 500 PSO iterations at 30,030 calls.

\paragraph*{Horns Rev~1 16-turbine block} We use the Horns Rev~1 offshore site from PyWake~\cite{PyWake}: Vestas V80 turbines (rotor diameter 80~m, 2~MW) with tabulated power and thrust curves and a 25~m/s cut-out, the measured 12-sector wind climate, and the parallelogram spanned by the installed turbines as boundary. We use the Jensen wake with $K=0.04$, $5^\circ$ direction and 1~m/s speed bins (agreement with PyWake's Jensen model within 0.3\% in AEP), and optimize PyWake's 16-turbine example block with $4D$ (320~m) minimum spacing.

\paragraph*{IEA37 Case Study~1} The optimization-only Case Study~1 of IEA Wind Task~37~\cite{Baker2019,IEA37repo} fixes the wake model and leaves the optimizer free. We use its 16- and 36-turbine scenarios (circular boundaries of radius 1300 and 2000~m, minimum spacing $2D=260$~m, IEA 3.35-MW reference turbine with $D=130$~m, 16 wind directions with a constant speed of 9.8~m/s, simplified Bastankhah Gaussian wake). Our implementation reproduces the official AEP calculator to a relative deviation below $10^{-11}$. All methods are run with the penalty~\eqref{eq:penalty} applied to the AEP loss, at 6,030 and 30,030 calls with 30 seeds, and compared with the submitted layouts~\cite{Baker2019,IEA37repo}.

\subsection{Performance Measures and Statistical Analysis}
For each run we record the final objective, the wake loss (as a percentage of the wake-free objective), whether the final layout is feasible (spacing and boundary satisfied within $10^{-6}$~m), the convergence curve (best feasible objective at 201 equally spaced checkpoints), and the wall-clock time. Means and standard deviations are computed over the feasible runs. Within each case, the 30 seed-paired runs are compared by two-sided Wilcoxon signed-rank tests of PSO-VNS against each other method, with Holm's correction~\cite{Holm1979} over the comparisons of the case and matched-pairs rank-biserial correlations $r_{\rm rb}$ as effect sizes; a feasible run beats an infeasible one. Across the 68 cases, the methods are ranked within each case by their mean feasible objective and compared by a Friedman test with the Iman--Davenport correction and Holm-adjusted post hoc tests on the average ranks~\cite{Demsar2006}. Because mean-rank post hoc tests depend on the set of compared methods~\cite{Benavoli2016}, we also compare PSO-VNS with each method by Wilcoxon signed-rank tests on the per-case mean wake losses.

\subsection{Implementation Verification and Environment}
\label{sec:validation}
Our vectorized objective agrees with the original implementation within a relative difference of $10^{-6}$ and reproduces the objective values of 720 archived final layouts within $8.7\times10^{-11}$; re-running the original SSA, LX-SSA, PSO and DE code with the archived settings reproduces the archived result distributions (22 of 24 Mann--Whitney comparisons with $p\ge0.05$). All runs used a Linux container (\TBD{CPU model and clock}, 4 virtual cores) with Python~3.11.15, NumPy~2.4.6 and SciPy~1.17.1.

\section{Results on the Benchmark}
\label{sec:results}
This section reports the 68 benchmark cases (\TBD{number of runs} runs) and the ablation. Per-case means and standard deviations are given in Tables~S1--S6 of the supplementary material.

\subsection{Overall Comparison}
\label{sec:overall}
\begin{table*}[!t]
\centering
\caption{Case-Level Analysis over the 68 Benchmark Cases. Avg.\ Rank: Average Rank of the Mean Feasible Objective (1 = Best); ``Best'': Cases in Which the Method Alone Ranks First; ``Feas.'': Feasible Runs; ``$<$15'': Cases with Fewer Than 15 Feasible Runs (Ranked Last); $p_z$: Holm-Adjusted $p$ of the Average-Rank Test Against PSO-VNS; Because Mean-Rank Tests Depend on the Pool of Methods~\cite{Benavoli2016}, $p_W$, $r_{\rm rb}$ and $\overline{\Delta L}$ Give the Holm-Adjusted Wilcoxon Test on the Per-Case Mean Wake Losses, the Rank-Biserial Correlation (Positive = PSO-VNS Better) and the Mean Wake-Loss Difference (Percentage Points, PSO-VNS Minus Method)}
\label{tab:friedman68}
\scriptsize\setlength{\tabcolsep}{4pt}
\begin{tabular}{lcccccccc}
\toprule
Method & Avg.\ rank & Best & Feas.\ (\%) & $<$15 & $p_z$ & $p_W$ & $r_{\rm rb}$ & $\overline{\Delta L}$ (pp) \\
\midrule
PSO-VNS & \TBD{} & \TBD{} & \TBD{} & \TBD{} & -- & -- & -- & -- \\
PSO & \TBD{} & \TBD{} & \TBD{} & \TBD{} & \TBD{} & \TBD{} & \TBD{} & \TBD{} \\
SSA-VNS & \TBD{} & \TBD{} & \TBD{} & \TBD{} & \TBD{} & \TBD{} & \TBD{} & \TBD{} \\
SSA & \TBD{} & \TBD{} & \TBD{} & \TBD{} & \TBD{} & \TBD{} & \TBD{} & \TBD{} \\
LX-SSA & \TBD{} & \TBD{} & \TBD{} & \TBD{} & \TBD{} & \TBD{} & \TBD{} & \TBD{} \\
DE & \TBD{} & \TBD{} & \TBD{} & \TBD{} & \TBD{} & \TBD{} & \TBD{} & \TBD{} \\
VNS & \TBD{} & \TBD{} & \TBD{} & \TBD{} & \TBD{} & \TBD{} & \TBD{} & \TBD{} \\
MS-SLSQP & \TBD{} & \TBD{} & \TBD{} & \TBD{} & \TBD{} & \TBD{} & \TBD{} & \TBD{} \\
\bottomrule
\multicolumn{9}{l}{Friedman $\chi^2_F=$ \TBD{} (7 d.f.), $p=$ \TBD{}; Iman--Davenport $F_F=$ \TBD{}, $p=$ \TBD{}.}
\end{tabular}
\end{table*}
\begin{table}[!t]
\centering
\caption{Pairwise Outcome of PSO-VNS Against Each Method: Cases in Which PSO-VNS Is Significantly Better / Not Different / Worse (Wilcoxon Signed-Rank Test, 30 Seed-Paired Runs, Holm-Adjusted over the Seven Comparisons of Each Case, $\alpha=0.05$); Last Row: Median Rank-Biserial Correlation (Positive = PSO-VNS Better)}
\label{tab:wtl}
\scriptsize\setlength{\tabcolsep}{1.6pt}
\begin{tabular}{llccccccccc}
\toprule
DS & $r$ (m) & Cases & PSO & \begin{tabular}{@{}c@{}}SSA-\\VNS\end{tabular} & SSA & \begin{tabular}{@{}c@{}}LX-\\SSA\end{tabular} & DE & VNS & \begin{tabular}{@{}c@{}}MS-\\SLSQP\end{tabular} \\
\midrule
I & 500 & 9 & \TBD{} & \TBD{} & \TBD{} & \TBD{} & \TBD{} & \TBD{} & \TBD{} \\
I & 750 & 11 & \TBD{} & \TBD{} & \TBD{} & \TBD{} & \TBD{} & \TBD{} & \TBD{} \\
I & 1000 & 14 & \TBD{} & \TBD{} & \TBD{} & \TBD{} & \TBD{} & \TBD{} & \TBD{} \\
II & 500 & 9 & \TBD{} & \TBD{} & \TBD{} & \TBD{} & \TBD{} & \TBD{} & \TBD{} \\
II & 750 & 11 & \TBD{} & \TBD{} & \TBD{} & \TBD{} & \TBD{} & \TBD{} & \TBD{} \\
II & 1000 & 14 & \TBD{} & \TBD{} & \TBD{} & \TBD{} & \TBD{} & \TBD{} & \TBD{} \\
\midrule
All & & 68 & \TBD{} & \TBD{} & \TBD{} & \TBD{} & \TBD{} & \TBD{} & \TBD{} \\
$\tilde r_{\rm rb}$ & & & \TBD{} & \TBD{} & \TBD{} & \TBD{} & \TBD{} & \TBD{} & \TBD{} \\
\bottomrule
\end{tabular}
\end{table}
\begin{figure}[!t]
\centering
\pipefig{\columnwidth}{figures_mpce/avg_ranks.pdf}{4cm}
\caption{Average rank of the eight methods over the 68 benchmark cases (1 = best), with the critical difference of the Holm-adjusted post hoc test. \TBD{regenerate with PSO-VNS}}
\label{fig:ranks}
\end{figure}

Table~\ref{tab:friedman68} and Fig.~\ref{fig:ranks} summarize the comparison over all cases. The Friedman test \TBD{rejects/does not reject} the hypothesis that the eight methods perform equally ($\chi^2_F=$ \TBD{}, $p=$ \TBD{}). PSO-VNS has the best average rank (\TBD{}) and alone ranks first in \TBD{} cases, followed by \TBD{ordering of the other methods with ranks}. In the Holm-adjusted post hoc tests, PSO-VNS ranks significantly better than \TBD{methods and $p_{\rm Holm}$}, and the Wilcoxon tests on the per-case means \TBD{agree / differ, esp.\ for PSO}.

The run-level tests in Table~\ref{tab:wtl} show where these differences come from. Against PSO, which forms its own first phase, PSO-VNS is significantly better in \TBD{} cases and worse in \TBD{} ($\tilde r_{\rm rb}=$ \TBD{}); \TBD{where the ties occur, e.g., small $N$}. Against SSA-VNS, the same design with a salp swarm first phase, it is better in \TBD{} and worse in \TBD{} cases. Against SSA, LX-SSA and DE, it is significantly better in \TBD{}, \TBD{} and \TBD{} cases, respectively. The comparison with VNS \TBD{W/T/L and interpretation}, and MS-SLSQP \TBD{outcome, incl.\ cases it wins}.

\subsection{Solution Quality, Feasibility and Convergence}
\begin{figure*}[!t]
\centering
\pipefig{0.8\textwidth}{figures_mpce/wakeloss_vs_n.pdf}{7cm}
\caption{Mean wake loss (percentage of the wake-free objective) of the feasible runs versus the prescribed number of turbines; a point is omitted when a method has no feasible run. \TBD{regenerate with PSO-VNS}}
\label{fig:wakeloss}
\end{figure*}
\begin{figure*}[!t]
\centering
\pipefig{0.8\textwidth}{figures_mpce/convergence_max.pdf}{7cm}
\caption{Convergence for the largest turbine count of each farm: median (line) and interquartile range (band) of the best feasible wake loss over the 30 runs versus the number of objective calls (logarithmic scale). A curve starts once at least half of the runs have found a feasible layout; PSO-VNS and PSO coincide up to 3,030 calls, where PSO-VNS switches to VNS. \TBD{regenerate with PSO-VNS}}
\label{fig:conv-max}
\end{figure*}

Fig.~\ref{fig:wakeloss} shows how the wake loss depends on the case. It grows with $N$, falls as the farm becomes larger, and is higher for the broader wind rose of Data Set~II. For $N\le3$, all methods stay within \TBD{}\% of the wake-free objective. For 15 turbines in the 1000-m farm (Data Set~I), the mean wake loss is \TBD{}\% for PSO-VNS, \TBD{}\% for PSO, \TBD{}\% for SSA-VNS, \TBD{}\% for VNS and \TBD{}\% for SSA; for 12 turbines in the 750-m farm (Data Set~II), the corresponding values are \TBD{five values}. In the largest case of every farm, PSO-VNS reduces the mean wake loss of PSO by \TBD{range} percentage points.

PSO-VNS ends with a feasible layout in \TBD{}\% of all runs; the other methods reach \TBD{} (PSO), \TBD{} (SSA-VNS), \TBD{} (SSA), \TBD{} (LX-SSA), \TBD{} (DE), \TBD{} (VNS) and \TBD{}\% (MS-SLSQP). The feasibility counts of the penalty-based methods are the same for both wind data sets, because as long as every candidate is infeasible, the penalty dominates $F_p$. \TBD{feasibility of the PSO phase at the switch in the densest cases (500~m, $N=10$) and after VNS}.

Fig.~\ref{fig:conv-max} shows the effect of the second phase directly. Up to 3,030 calls, PSO-VNS and PSO follow the same trajectory. After the switch, \TBD{PSO-VNS drops / PSO stagnates or keeps improving}. Over all feasible runs, the VNS phase removes on average \TBD{}\% (median \TBD{}\%) of the wake loss left at the end of the PSO phase, and it turns \TBD{} runs that were still infeasible at the switch into feasible ones. \TBD{comparison with the SSA-VNS, VNS and MS-SLSQP curves}.

\subsection{Effect of the PSO Coefficients}
\label{sec:psosetting}
The earlier version of this study ran PSO with $w=0.7$ and $c_1=c_2=2$, which violates \eqref{eq:poli}, on the same 68 cases with the same budget and seeds. With this setting, PSO had an average rank of \TBD{} among the methods of that comparison, returned a feasible layout in \TBD{}\% of the runs, and was significantly worse than the constriction PSO in \TBD{} of the 68 cases (mean wake loss \TBD{} percentage points higher). The constriction setting thus moves PSO from \TBD{rank position} to \TBD{rank position}, ahead of SSA, LX-SSA and SSA-VNS. We draw two conclusions. First, a PSO baseline outside the convergence region understates what PSO achieves on WFLOP, and comparisons that rely on it can overstate the advantage of a new method; we recommend checking \eqref{eq:poli}, or using the constriction coefficients \eqref{eq:constriction}, whenever PSO serves as a baseline. Second, the salp swarm methods, including our own LX-SSA, do not outperform a correctly configured PSO on this benchmark, which is why PSO rather than SSA forms the first phase of the proposed method.

\subsection{Ablation}
\label{sec:ablation}
\begin{table}[!t]
\centering
\caption{Ablation over the 68 Benchmark Cases (30 Seed-Paired Runs, 6,030 Calls). Top: Average Rank Among the Eight Variants and Feasible Runs. Bottom: Planned Contrasts, Cases in Which the First Variant Is Significantly Better / Not Different / Worse (Wilcoxon, Holm-Adjusted over the Contrasts of Each Case, $\alpha=0.05$), and Mean Difference of the Wake Loss $\overline{\Delta L}$ (Percentage Points; Negative = First Variant Better)}
\label{tab:ablation}
\scriptsize\setlength{\tabcolsep}{2.5pt}
\begin{tabular}{lcccc}
\toprule
Variant & Phase 1 & Phase 2 & Avg.\ rank & Feas.\ (\%) \\
\midrule
PSO-VNS & PSO & VNS & \TBD{} & \TBD{} \\
SSA-VNS & SSA & VNS & \TBD{} & \TBD{} \\
LX-SSA-VNS & LX-SSA & VNS & \TBD{} & \TBD{} \\
RS-VNS & random sampling & VNS & \TBD{} & \TBD{} \\
VNS & -- & VNS & \TBD{} & \TBD{} \\
PSO & PSO & -- & \TBD{} & \TBD{} \\
SSA & SSA & -- & \TBD{} & \TBD{} \\
LX-SSA & LX-SSA & -- & \TBD{} & \TBD{} \\
\bottomrule
\end{tabular}

\smallskip
\begin{tabular}{l>{\raggedright\arraybackslash}p{2.9cm}cc}
\toprule
Contrast & Isolates & W/T/L & $\overline{\Delta L}$ \\
\midrule
PSO-VNS vs.\ PSO & VNS phase & \TBD{} & \TBD{} \\
PSO-VNS vs.\ VNS & swarm start vs.\ best initial point & \TBD{} & \TBD{} \\
PSO-VNS vs.\ RS-VNS & swarm vs.\ random sampling & \TBD{} & \TBD{} \\
PSO-VNS vs.\ SSA-VNS & PSO vs.\ SSA as Phase~1 & \TBD{} & \TBD{} \\
PSO-VNS vs.\ LX-SSA-VNS & PSO vs.\ LX-SSA as Phase~1 & \TBD{} & \TBD{} \\
SSA-VNS vs.\ RS-VNS & SSA vs.\ random sampling & \TBD{} & \TBD{} \\
SSA-VNS vs.\ LX-SSA-VNS & Laplace step in the hybrid & \TBD{} & \TBD{} \\
LX-SSA vs.\ SSA & Laplace step alone & \TBD{} & \TBD{} \\
\bottomrule
\end{tabular}
\end{table}
\begin{figure*}[!t]
\centering
\pipefig{0.8\textwidth}{figures_mpce/ablation_convergence.pdf}{7cm}
\caption{Ablation, largest turbine count of each farm: median best feasible wake loss over 30 runs versus objective calls for PSO-VNS, PSO, VNS, RS-VNS, SSA-VNS and LX-SSA-VNS (dotted line: switch to VNS at 3,030 calls). \TBD{regenerate with the PSO-VNS ablation variants}}
\label{fig:ablation}
\end{figure*}

To see where the improvement of PSO-VNS comes from, we compare it with variants that each remove or replace one ingredient (Table~\ref{tab:ablation} and Fig.~\ref{fig:ablation}). The Friedman test over the eight variants gives $\chi^2_F=$ \TBD{} (7 d.f., $p=$ \TBD{}), and PSO-VNS has the best average rank (\TBD{}).

\emph{VNS phase.} Switching to VNS after 3,030 calls instead of continuing the swarm \TBD{effect vs.\ PSO: W/T/L and mean difference}. Since PSO-VNS and PSO share their first 3,030 calls, this contrast measures only the use of the second half of the budget.

\emph{Swarm start.} Starting VNS from the PSO global best rather than from the best initial point \TBD{effect vs.\ VNS}. The random-multistart control RS-VNS spends the same 3,015 calls of Phase~1 on independent random layouts; \TBD{effect vs.\ RS-VNS: does the PSO phase add value beyond a good random start?}. For SSA the answer is negative: SSA-VNS and RS-VNS are statistically tied in \TBD{} of the 68 cases.

\emph{Choice of the Phase-1 swarm.} With the same VNS phase, PSO-VNS \TBD{W/T/L} against SSA-VNS and \TBD{W/T/L} against LX-SSA-VNS, so a convergent PSO is \TBD{the better / an equally good} first phase. \TBD{brief interpretation, e.g., feasibility or quality of the start at the switch}.

\emph{Laplace step.} The Laplace step of LX-SSA gave no gain: without the VNS phase, LX-SSA is \TBD{W/T/L} against SSA, and inside the hybrid, SSA-VNS is \TBD{W/T/L} against LX-SSA-VNS (mean wake-loss difference \TBD{} percentage points).

\section{Results Beyond the Benchmark}
\label{sec:beyond}

\subsection{Feasibility-Preserving Initialization and Budget Scaling}
\label{sec:feasinit}
\label{sec:budget}
\begin{table}[!t]
\centering
\caption{Six Largest Benchmark Cases: Mean Wake Loss (\%) and Feasible Runs (\%) at 6,030 Calls with Random and Feasibility-Preserving Initialization, and Average Rank at 6,030, 30,030 and 120,030 Calls (Random Initialization)}
\label{tab:feasbudget}
\scriptsize\setlength{\tabcolsep}{2.4pt}
\begin{tabular}{lccccccc}
\toprule
& \multicolumn{2}{c}{Random init.} & \multicolumn{2}{c}{Feasible init.} & \multicolumn{3}{c}{Avg.\ rank} \\
\cmidrule(lr){2-3}\cmidrule(lr){4-5}\cmidrule(lr){6-8}
Method & Loss & Feas. & Loss & Feas. & 6,030 & 30,030 & 120,030 \\
\midrule
PSO-VNS & \TBD{} & \TBD{} & \TBD{} & \TBD{} & \TBD{} & \TBD{} & \TBD{} \\
SSA-VNS & \TBD{} & \TBD{} & \TBD{} & \TBD{} & \TBD{} & \TBD{} & \TBD{} \\
LX-SSA-VNS & \TBD{} & \TBD{} & \TBD{} & \TBD{} & \TBD{} & \TBD{} & \TBD{} \\
RS-VNS & \TBD{} & \TBD{} & \TBD{} & \TBD{} & \TBD{} & \TBD{} & \TBD{} \\
VNS & \TBD{} & \TBD{} & \TBD{} & \TBD{} & \TBD{} & \TBD{} & \TBD{} \\
PSO & \TBD{} & \TBD{} & \TBD{} & \TBD{} & \TBD{} & \TBD{} & \TBD{} \\
SSA & \TBD{} & \TBD{} & \TBD{} & \TBD{} & \TBD{} & \TBD{} & \TBD{} \\
LX-SSA & \TBD{} & \TBD{} & \TBD{} & \TBD{} & \TBD{} & \TBD{} & \TBD{} \\
DE & \TBD{} & \TBD{} & \TBD{} & \TBD{} & \TBD{} & \TBD{} & \TBD{} \\
MS-SLSQP & \TBD{} & \TBD{} & \TBD{} & \TBD{} & \TBD{} & \TBD{} & \TBD{} \\
\bottomrule
\end{tabular}
\end{table}
\begin{figure*}[!t]
\centering
\pipefig{0.8\textwidth}{figures_mpce/budget_scaling.pdf}{8cm}
\caption{Budget scaling: mean wake loss of the six largest benchmark cases and AEP loss of the Horns Rev~1 block versus the budget (6,030, 30,030 and 120,030 calls, random initialization). \TBD{regenerate with PSO-VNS; check panel description}}
\label{fig:budget}
\end{figure*}

Table~\ref{tab:feasbudget} summarizes the six largest benchmark cases. With feasibility-preserving initialization, \TBD{effect on feasibility rates, in particular for PSO, PSO-VNS and DE in the dense cases; effect on wake loss and on the ranking}. With 5 and 20 times the budget (Fig.~\ref{fig:budget}), \TBD{does the ranking change; does the advantage of PSO-VNS over PSO, VNS, RS-VNS and MS-SLSQP grow or shrink}.

\subsection{Horns Rev~1 16-Turbine Block}
\label{sec:hornsrev-site}
\begin{table}[!t]
\centering
\caption{Horns Rev~1 16-Turbine Block (Wake-Free AEP 149.26 GWh/yr): Mean AEP (GWh/yr) over the Feasible Runs of 30 Seeds at 6,030 Calls, Feasible Runs, Holm-Adjusted Wilcoxon $p$ of PSO-VNS vs.\ Each Method, and Mean AEP Loss (\%) at 6,030 Calls with Random (R) and Feasibility-Preserving (F) Initialization and at 30,030 and 120,030 Calls (R; 10 Seeds at 120,030)}
\label{tab:hr-site}
\scriptsize\setlength{\tabcolsep}{2.2pt}
\begin{tabular}{lccccccc}
\toprule
& \multicolumn{3}{c}{6,030 calls, R} & \multicolumn{4}{c}{Loss (\%)} \\
\cmidrule(lr){2-4}\cmidrule(lr){5-8}
Layout / method & AEP & Feas. & $p_{\rm Holm}$ & 6k R & 6k F & 30k & 120k \\
\midrule
Installed layout & 139.51 & -- & -- & 6.53 & -- & -- & -- \\
\midrule
PSO-VNS & \TBD{} & \TBD{} & -- & \TBD{} & \TBD{} & \TBD{} & \TBD{} \\
SSA-VNS & \TBD{} & \TBD{} & \TBD{} & \TBD{} & \TBD{} & \TBD{} & \TBD{} \\
LX-SSA-VNS & \TBD{} & \TBD{} & \TBD{} & \TBD{} & \TBD{} & \TBD{} & \TBD{} \\
RS-VNS & \TBD{} & \TBD{} & \TBD{} & \TBD{} & \TBD{} & \TBD{} & \TBD{} \\
VNS & \TBD{} & \TBD{} & \TBD{} & \TBD{} & \TBD{} & \TBD{} & \TBD{} \\
PSO & \TBD{} & \TBD{} & \TBD{} & \TBD{} & \TBD{} & \TBD{} & \TBD{} \\
SSA & \TBD{} & \TBD{} & \TBD{} & \TBD{} & \TBD{} & \TBD{} & \TBD{} \\
LX-SSA & \TBD{} & \TBD{} & \TBD{} & \TBD{} & \TBD{} & \TBD{} & \TBD{} \\
DE & \TBD{} & \TBD{} & \TBD{} & \TBD{} & \TBD{} & \TBD{} & \TBD{} \\
MS-SLSQP & \TBD{} & \TBD{} & \TBD{} & \TBD{} & \TBD{} & \TBD{} & \TBD{} \\
\bottomrule
\end{tabular}
\end{table}

Table~\ref{tab:hr-site} gives the results for the 16-turbine block; the installed layout, a regular $7D$ grid, reaches 139.51~GWh/yr (wake loss 6.53\%). PSO-VNS attains \TBD{mean AEP, feasibility and significance against the other methods}. \TBD{whether PSO finds feasible layouts inside the parallelogram from random starts, and how feasible initialization changes this}. \TBD{whether any method improves on the installed layout at 6,030 calls and at the larger budgets}. The best layouts are shown in Fig.~\ref{fig:layouts}.

\subsection{IEA37 Case Study~1}
\label{sec:iea37}
\begin{table}[!t]
\centering
\caption{IEA37 Case Study~1, 16- and 36-Turbine Scenarios: AEP in MWh of the Best Run of Each Method at 6,030 and 30,030 Calls, Compared with the Example Layout and the Best Feasible Published Layout~\cite{Baker2019,IEA37repo}}
\label{tab:iea37}
\scriptsize\setlength{\tabcolsep}{2.5pt}
\begin{tabular}{lcccc}
\toprule
& \multicolumn{2}{c}{16 turbines ($r=1300$~m)} & \multicolumn{2}{c}{36 turbines ($r=2000$~m)} \\
\cmidrule(lr){2-3}\cmidrule(lr){4-5}
Method / source & 6,030 & 30,030 & 6,030 & 30,030 \\
\midrule
Example layout & \multicolumn{2}{c}{366,941.6} & \multicolumn{2}{c}{737,883.1} \\
Best feasible published$^{a}$ & \multicolumn{2}{c}{418,924.4} & \multicolumn{2}{c}{863,676.3} \\
\midrule
PSO-VNS & \TBD{} & \TBD{} & \TBD{} & \TBD{} \\
PSO & \TBD{} & \TBD{} & \TBD{} & \TBD{} \\
SSA-VNS & \TBD{} & \TBD{} & \TBD{} & \TBD{} \\
SSA & \TBD{} & \TBD{} & \TBD{} & \TBD{} \\
LX-SSA & \TBD{} & \TBD{} & \TBD{} & \TBD{} \\
DE & \TBD{} & \TBD{} & \TBD{} & \TBD{} \\
VNS & \TBD{} & \TBD{} & \TBD{} & \TBD{} \\
MS-SLSQP & \TBD{} & \TBD{} & \TBD{} & \TBD{} \\
\bottomrule
\multicolumn{5}{p{0.92\columnwidth}}{$^{a}$Participant~4 in both scenarios (SNOPT with wake expansion continuation \TBD{verify against Baker et al.\ (2019)}). A higher submitted AEP places turbines up to 3.5~m outside the boundary and is not counted.}
\end{tabular}
\end{table}
\begin{figure*}[!t]
\centering
\pipefig{0.75\textwidth}{figures_mpce/iea37_layouts.pdf}{5cm}\\[2pt]
\pipefig{0.75\textwidth}{figures_mpce/hr16.pdf}{5.4cm}
\caption{Top: best PSO-VNS layouts for the IEA37 16- and 36-turbine scenarios at 30,030 calls, with the best feasible published layouts for comparison. Bottom: Horns Rev~1 16-turbine block, installed layout and best PSO-VNS layout. \TBD{regenerate with PSO-VNS; check panel content of both files}}
\label{fig:layouts}
\end{figure*}

Table~\ref{tab:iea37} and Fig.~\ref{fig:layouts} place our methods in the context of IEA37 Case Study~1, in which the participants chose their own optimizers and budgets~\cite{Baker2019}. Starting from the example layouts, the best feasible submitted layouts reach 418,924.4~MWh for 16 turbines and 863,676.3~MWh for 36 turbines. \TBD{rank of the best and mean PSO-VNS layouts among the submitted layouts; gap to the best feasible published layout at 6,030 and 30,030 calls}. Since most participants used gradient-based methods and budgets that are not reported in a comparable form, this comparison is indicative only.

\subsection{Published Kusiak--Song Results and Computational Cost}
\label{sec:literature}
\label{sec:runtime}
The Kusiak--Song benchmark has been solved with evolution strategies~\cite{Kusiak2010}, ant colony optimization~\cite{Eroglu2012}, particle filtering~\cite{Eroglu2013} and biogeography-based optimization~\cite{Bansal2017}; a comparison with the published values is given in the supplementary material. \TBD{one-sentence summary, pending author verification of the PDFs}. At equal numbers of objective calls, all methods need almost the same time (\TBD{range of ms per call}); the objective evaluation dominates the cost, and one 6,030-call run of PSO-VNS takes \TBD{range}~s, depending on $N$.

\section{Sensitivity to the Modeling Assumptions}
\label{sec:robustness}
\subsection{Power Curve and Wake Model}
\label{sec:powercurve}
\label{sec:gaussian}
\begin{table}[!t]
\centering
\caption{Robustness to the Benchmark Model: All Feasible Final Layouts Re-Evaluated (Not Re-Optimized). $\Delta$: Mean Relative Change of the Objective (\%); Rank: Average Rank of PSO-VNS over the 68 Cases; Best: Best-Ranked Method; $\bar\tau$: Mean Kendall Correlation Between the Benchmark and Alternative Orderings Within a Case; Same: Cases (\%) with Unchanged Best Method. The Ranks of All Methods Are Given in the Supplementary Material}
\label{tab:robust-final}
\scriptsize\setlength{\tabcolsep}{2.4pt}
\begin{tabular}{lccccccc}
\toprule
& \multicolumn{2}{c}{$\Delta$ (\%)} & & & & \\
\cmidrule(lr){2-3}
Model & DS I & DS II & Rank & Best & $\bar\tau$ & Same \\
\midrule
Benchmark & +0.0 & +0.0 & \TBD{} & \TBD{} & -- & -- \\
Cubic curve & \TBD{} & \TBD{} & \TBD{} & \TBD{} & \TBD{} & \TBD{} \\
Cubic, 25 m/s cut-out & \TBD{} & \TBD{} & \TBD{} & \TBD{} & \TBD{} & \TBD{} \\
Gaussian wake & \TBD{} & \TBD{} & \TBD{} & \TBD{} & \TBD{} & \TBD{} \\
\bottomrule
\end{tabular}
\end{table}

We re-evaluated every feasible final layout of the main comparison, without re-optimization, with 1) a cubic power curve~\cite{Carrillo2013}, $f_{\rm c}(s)=P_{\text{rated}}(s^3-s_{\text{cut-in}}^3)/(s_{\text{rated}}^3-s_{\text{cut-in}}^3)$ between the cut-in and rated speeds, 2) the same curve with a 25~m/s cut-out, and 3) the Gaussian wake model of Bastankhah and Port\'e-Agel~\cite{Bastankhah2014} with growth rate $K_{\rm G}=0.04$ and the same superposition and Weibull scaling as the benchmark (equations in the supplementary material).

Table~\ref{tab:robust-final} shows two kinds of sensitivity. First, the absolute energy depends strongly on the model: the linear benchmark curve overstates the expected power relative to the cubic curve by \TBD{}\% for Data Set~I and \TBD{}\% for Data Set~II, so the energy values in this paper are benchmark values and not AEP predictions for a specific turbine. Second, the ranking of the algorithms is \TBD{stable/unstable} under the cubic curve ($\bar\tau=$ \TBD{}) and \TBD{} under the Gaussian model ($\bar\tau=$ \TBD{}; same best method in \TBD{}\% of the cases), because layouts optimized for the sharp Jensen wake cone are judged differently by a smooth Gaussian deficit. \TBD{rank of PSO-VNS under the Gaussian model}.

\subsection{Minimum Spacing and Boundary Handling}
\label{sec:spacing}
\label{sec:boundary}
The $4D$ spacing is inherited from the benchmark literature. A multi-start packing search constructed feasible layouts with up to 13, 26 and 42 turbines at $4D$, 8, 19 and 29 at $5D$, and 7, 13 and 21 at $6D$ for $r=500$, 750 and 1000~m (constructive lower bounds), so the 500-m cases with $N\ge9$ (at $5D$) or $N\ge8$ (at $6D$) would be infeasible. The optimized layouts often lie on the $4D$ constraint: of the \TBD{} feasible final layouts of the eight methods, \TBD{}\% also satisfy $5D$ and \TBD{}\% satisfy $6D$, and for $N=11$--15 these shares fall to \TBD{}\% and \TBD{}\%, so layouts optimized for $4D$ cannot be transferred to a site with larger spacing requirements without re-optimization. Re-optimizing six representative cases with LX-SSA and SSA at $5D$ and $6D$ (Table~S7 of the supplementary material) lowers the mean objective by up to 0.95\% but does not reverse the comparison between the two methods.

The boundary rule also matters: an otherwise identical SSA that projects every turbine radially onto the circle instead of clipping the coordinates attains higher mean objectives than the original SSA in six representative cases at 3,030 calls, by 26 to 649 objective units (Mann--Whitney $p<10^{-4}$ in every case). We therefore keep the rule of Section~\ref{sec:constraints} fixed for all methods.

\section{Limitations}
\label{sec:limits}
Our results should be read within the following limits. The components of PSO-VNS are standard, and its gain comes mainly from \TBD{VNS phase / swarm start, according to the ablation}; in single cases, PSO-VNS differs significantly from PSO in \TBD{} and from VNS in \TBD{} of the 68 cases. The budget study covers up to 120,030 calls but only the six largest benchmark cases and the Horns Rev~1 block, and the largest farms have 36 turbines; farms with 80 or more turbines, where penalty-based search struggles to find feasible layouts, are outside our scope. The PSO coefficients are the standard constriction values, and the VNS settings and $\omega=0.5$ are untuned defaults. Our statement on the setting $w=0.7$, $c_1=c_2=2$ rests on the convergence theory and our own benchmark; we did not re-run the comparisons of other studies. The absolute energy depends strongly on the Jensen model, the linear power curve, the $4D$ spacing and the boundary rule, \TBD{the ranking under the Gaussian model}, and we did not re-optimize with the alternative models. Finally, $N$ is prescribed, and terrain, mixed turbine types, cables, grid connection, noise, environmental constraints and economic objectives are outside the scope of this study.

\section{Conclusion}
\label{sec:conclusion}
This work presents PSO-VNS, a two-phase algorithm for the continuous wind farm layout problem in which PSO with constriction coefficients explores the layout space for the first half of the budget and the basic VNS then refines the global best layout. We evaluated it against seven methods on 68 benchmark cases at equal budgets, in an ablation, with feasibility-preserving initialization and larger budgets, on a Horns Rev~1 block and on IEA37 Case Study~1.

PSO-VNS obtained the best average rank (\TBD{}), was \TBD{significance summary} and returned a feasible layout in \TBD{}\% of the runs. The ablation showed that \TBD{contributions of the VNS phase and of the swarm start, incl.\ vs.\ RS-VNS}, that PSO is a better first phase than SSA or LX-SSA, and that the Laplace step of our earlier LX-SSA did not improve the results, either alone or inside the hybrid. PSO configured inside its convergence region outperformed the salp swarm methods, whereas the frequently used setting $w=0.7$, $c_1=c_2=2$ violates the convergence condition and understates PSO; baselines in WFLOP comparisons should be checked against this condition. \TBD{one sentence each on feasible initialization, budget scaling, Horns Rev~1 and IEA37}. The ranking of the methods was \TBD{} under a cubic power curve and \TBD{} under a Gaussian wake model.

In future work, we will study a switch between the two phases that is triggered by the stagnation of the swarm, investigate which properties of the swarm phase make it a good starting point for VNS, and extend the approach to larger farms and higher-fidelity wake and power models.

\section*{Acknowledgment}
\TBD{acknowledgment, or delete this section}

\section*{Data Availability}
The source code, the per-run results (including the final coordinates and convergence curves of all runs) and the scripts that generate all tables and figures are openly available at \TBD{Zenodo DOI}. The supplementary material contains the schematics of the farm, the wake model and the wake cone (Figs.~S1--S3), the pseudocode of SSA and LX-SSA, the per-case results of the 68 benchmark cases (Tables~S1--S6), the minimum-spacing re-optimization (Table~S7), the budget-split study of SSA-VNS, the comparison with published Kusiak--Song results, the coordinates of the best layouts, and further convergence, box-plot, feasibility and layout figures. \TBD{assemble supplementary file; check S-numbering}

\begin{thebibliography}{99}

\bibitem{Mosetti1994}
G. Mosetti, C. Poloni, and B. Diviacco, ``Optimization of wind turbine positioning in large windfarms by means of a genetic algorithm,'' \textit{J. Wind Eng. Ind. Aerodyn.}, vol. 51, no. 1, pp. 105--116, Jan. 1994, doi: 10.1016/0167-6105(94)90080-9.

\bibitem{Grady2005}
S. A. Grady, M. Y. Hussaini, and M. M. Abdullah, ``Placement of wind turbines using genetic algorithms,'' \textit{Renew. Energy}, vol. 30, no. 2, pp. 259--270, Feb. 2005, doi: 10.1016/j.renene.2004.05.007.

\bibitem{Kusiak2010}
A. Kusiak and Z. Song, ``Design of wind farm layout for maximum wind energy capture,'' \textit{Renew. Energy}, vol. 35, no. 3, pp. 685--694, Mar. 2010, doi: 10.1016/j.renene.2009.08.019.

\bibitem{Eroglu2012}
Y. Ero\u{g}lu and S. U. Se\c{c}kiner, ``Design of wind farm layout using ant colony algorithm,'' \textit{Renew. Energy}, vol. 44, pp. 53--62, Aug. 2012, doi: 10.1016/j.renene.2011.12.013.

\bibitem{Eroglu2013}
Y. Ero\u{g}lu and S. U. Se\c{c}kiner, ``Wind farm layout optimization using particle filtering approach,'' \textit{Renew. Energy}, vol. 58, pp. 95--107, Oct. 2013, doi: 10.1016/j.renene.2013.02.019.

\bibitem{Bansal2017}
J. C. Bansal and P. Farswan, ``Wind farm layout using biogeography based optimization,'' \textit{Renew. Energy}, vol. 107, pp. 386--402, Jul. 2017, doi: 10.1016/j.renene.2017.01.064.

\bibitem{Elkinton2008}
C. N. Elkinton, J. F. Manwell, and J. G. McGowan, ``Algorithms for offshore wind farm layout optimization,'' \textit{Wind Eng.}, vol. 32, no. 1, pp. 67--84, Jan. 2008, doi: 10.1260/030952408784305877.

\bibitem{Asaah2021}
P. Asaah, L. Hao, and J. Ji, ``Optimal placement of wind turbines in wind farm layout using particle swarm optimization,'' \textit{J. Mod. Power Syst. Clean Energy}, vol. 9, no. 2, pp. 367--375, Mar. 2021, doi: 10.35833/MPCE.2019.000087.

\bibitem{Srikakulapu2018}
R. Srikakulapu and U. Vinatha, ``Optimized design of collector topology for offshore wind farm based on ant colony optimization with multiple travelling salesman problem,'' \textit{J. Mod. Power Syst. Clean Energy}, vol. 6, no. 6, pp. 1181--1192, Nov. 2018, doi: 10.1007/s40565-018-0386-4.

\bibitem{Hou2019}
P. Hou, J. Zhu, K. Ma, G. Yang, W. Hu, and Z. Chen, ``A review of offshore wind farm layout optimization and electrical system design methods,'' \textit{J. Mod. Power Syst. Clean Energy}, vol. 7, no. 5, pp. 975--986, Sep. 2019, doi: 10.1007/s40565-019-0550-5.

\bibitem{Ju2019}
X. Ju and F. Liu, ``Wind farm layout optimization using self-informed genetic algorithm with information guided exploitation,'' \textit{Appl. Energy}, vol. 248, pp. 429--445, Aug. 2019, doi: 10.1016/j.apenergy.2019.04.084.

\bibitem{Bai2022}
F. Bai, X. Ju, S. Wang, W. Zhou, and F. Liu, ``Wind farm layout optimization using adaptive evolutionary algorithm with Monte Carlo tree search reinforcement learning,'' \textit{Energy Convers. Manage.}, vol. 252, Art. no. 115047, Jan. 2022, doi: 10.1016/j.enconman.2021.115047.

\bibitem{Nagpal2021}
S. V. Nagpal, M. V. Liu, and C. L. Anderson, ``A comparison of deterministic refinement techniques for wind farm layout optimization,'' \textit{Renew. Energy}, vol. 168, pp. 581--592, May 2021, doi: 10.1016/j.renene.2020.12.043.

\bibitem{Cazzaro2022}
D. Cazzaro and D. Pisinger, ``Variable neighborhood search for large offshore wind farm layout optimization,'' \textit{Comput. Oper. Res.}, vol. 138, Art. no. 105588, Feb. 2022, doi: 10.1016/j.cor.2021.105588.

\bibitem{Quaeghebeur2021}
E. Quaeghebeur, R. Bos, and M. B. Zaaijer, ``Wind farm layout optimization using pseudo-gradients,'' \textit{Wind Energy Sci.}, vol. 6, no. 3, pp. 815--839, Jun. 2021, doi: 10.5194/wes-6-815-2021.

\bibitem{Tao2020}
S. Tao, Q. Xu, A. Feij\'oo, G. Zheng, and J. Zhou, ``Wind farm layout optimization with a three-dimensional Gaussian wake model,'' \textit{Renew. Energy}, vol. 159, pp. 553--569, Oct. 2020, doi: 10.1016/j.renene.2020.06.003.

\bibitem{Yang2023}
K. Yang, X. Deng, Z. Ti, S. Yang, S. Huang, and Y. Wang, ``A data-driven layout optimization framework of large-scale wind farms based on machine learning,'' \textit{Renew. Energy}, vol. 218, Art. no. 119240, Dec. 2023, doi: 10.1016/j.renene.2023.119240.

\bibitem{Wang2024}
Z. Wang, Y. Tu, K. Zhang, Z. Han, Y. Cao, and D. Zhou, ``An optimization framework for wind farm layout design using CFD-based Kriging model,'' \textit{Ocean Eng.}, vol. 293, Art. no. 116644, Feb. 2024, doi: 10.1016/j.oceaneng.2023.116644.

\bibitem{Li2025}
P. Li, Y. Che, A. Hua, L. Wang, M. Zheng, and X. Guo, ``A data-physics hybrid-driven layout optimization framework for large-scale wind farms,'' \textit{Appl. Energy}, vol. 392, Art. no. 125908, 2025, doi: 10.1016/j.apenergy.2025.125908.

\bibitem{Baker2019}
N. F. Baker, A. P. J. Stanley, J. J. Thomas, A. Ning, and K. Dykes, ``Best practices for wake model and optimization algorithm selection in wind farm layout optimization,'' in \textit{Proc. AIAA SciTech Forum}, San Diego, CA, USA, Jan. 2019, Paper AIAA 2019-0540, doi: 10.2514/6.2019-0540.

\bibitem{Clerc2002}
M. Clerc and J. Kennedy, ``The particle swarm---Explosion, stability, and convergence in a multidimensional complex space,'' \textit{IEEE Trans. Evol. Comput.}, vol. 6, no. 1, pp. 58--73, Feb. 2002, doi: 10.1109/4235.985692.

\bibitem{Poli2009}
R. Poli, ``Mean and variance of the sampling distribution of particle swarm optimizers during stagnation,'' \textit{IEEE Trans. Evol. Comput.}, vol. 13, no. 4, pp. 712--721, Aug. 2009, doi: 10.1109/TEVC.2008.2011744.

\bibitem{Mirjalili2017}
S. Mirjalili, A. H. Gandomi, S. Z. Mirjalili, S. Saremi, H. Faris, and S. M. Mirjalili, ``Salp swarm algorithm: A bio-inspired optimizer for engineering design problems,'' \textit{Adv. Eng. Softw.}, vol. 114, pp. 163--191, Dec. 2017, doi: 10.1016/j.advengsoft.2017.07.002.

\bibitem{Solanki2023}
P. Solanki and K. Deep, ``Laplacian salp swarm algorithm for continuous optimization,'' \textit{Int. J. Syst. Assur. Eng. Manag.}, 2023, doi: 10.1007/s13198-023-01935-y.

\bibitem{Mladenovic1997}
N. Mladenovi\'c and P. Hansen, ``Variable neighborhood search,'' \textit{Comput. Oper. Res.}, vol. 24, no. 11, pp. 1097--1100, Nov. 1997, doi: 10.1016/S0305-0548(97)00031-2.

\bibitem{Hansen2001}
P. Hansen and N. Mladenovi\'c, ``Variable neighborhood search: Principles and applications,'' \textit{Eur. J. Oper. Res.}, vol. 130, no. 3, pp. 449--467, May 2001, doi: 10.1016/S0377-2217(00)00100-4.

\bibitem{Jensen1983}
N. O. Jensen, ``A note on wind generator interaction,'' Ris\o\ Nat. Lab., Roskilde, Denmark, Tech. Rep. Ris\o-M-2411, 1983.

\bibitem{Katic1986}
I. Kati\'c, J. H\o jstrup, and N. O. Jensen, ``A simple model for cluster efficiency,'' in \textit{Proc. Eur. Wind Energy Assoc. Conf. Exhib. (EWEC)}, Rome, Italy, Oct. 1986, vol. 1, pp. 407--410.

\bibitem{Kennedy1995}
J. Kennedy and R. Eberhart, ``Particle swarm optimization,'' in \textit{Proc. IEEE Int. Conf. Neural Netw. (ICNN)}, Perth, WA, Australia, Nov. 1995, vol. 4, pp. 1942--1948, doi: 10.1109/ICNN.1995.488968.

\bibitem{Shi1998}
Y. Shi and R. Eberhart, ``A modified particle swarm optimizer,'' in \textit{Proc. IEEE Int. Conf. Evol. Comput. (ICEC)}, Anchorage, AK, USA, May 1998, pp. 69--73, doi: 10.1109/ICEC.1998.699146.

\bibitem{Deep2007}
K. Deep and M. Thakur, ``A new crossover operator for real coded genetic algorithms,'' \textit{Appl. Math. Comput.}, vol. 188, no. 1, pp. 895--911, May 2007, doi: 10.1016/j.amc.2006.10.047.

\bibitem{Mladenovic2008}
N. Mladenovi\'c, M. Dra\v{z}i\'c, V. Kova\v{c}evi\'c-Vuj\v{c}i\'c, and M. \v{C}angalovi\'c, ``General variable neighborhood search for the continuous optimization,'' \textit{Eur. J. Oper. Res.}, vol. 191, no. 3, pp. 753--770, Dec. 2008, doi: 10.1016/j.ejor.2006.12.064.

\bibitem{Storn1997}
R. Storn and K. Price, ``Differential evolution---A simple and efficient heuristic for global optimization over continuous spaces,'' \textit{J. Global Optim.}, vol. 11, no. 4, pp. 341--359, Dec. 1997, doi: 10.1023/A:1008202821328.

\bibitem{Kraft1988}
D. Kraft, ``A software package for sequential quadratic programming,'' DFVLR, Inst. f\"ur Dynamik der Flugsysteme, Oberpfaffenhofen, Germany, Tech. Rep. DFVLR-FB 88-28, 1988.

\bibitem{PyWake}
DTU Wind Energy, ``PyWake: Open-source static wake modeling framework,'' Python package \texttt{py\_wake}, ver. 2.6.20 (Horns Rev~1 site data in \texttt{py\_wake/examples/data/hornsrev1.py}), MIT License. [Online]. Available: \url{https://gitlab.windenergy.dtu.dk/TOPFARM/PyWake}

\bibitem{IEA37repo}
N. F. Baker \textit{et al.}, ``IEA Wind Task 37 wind farm layout optimization case studies: Case files and submitted results.'' [Online]. Available: \url{https://github.com/IEAWindTask37/iea37-wflo-casestudies}

\bibitem{Holm1979}
S. Holm, ``A simple sequentially rejective multiple test procedure,'' \textit{Scand. J. Statist.}, vol. 6, no. 2, pp. 65--70, 1979.

\bibitem{Demsar2006}
J. Dem\v{s}ar, ``Statistical comparisons of classifiers over multiple data sets,'' \textit{J. Mach. Learn. Res.}, vol. 7, pp. 1--30, Jan. 2006.

\bibitem{Benavoli2016}
A. Benavoli, G. Corani, and F. Mangili, ``Should we really use post-hoc tests based on mean-ranks?'' \textit{J. Mach. Learn. Res.}, vol. 17, no. 5, pp. 1--10, 2016.

\bibitem{Carrillo2013}
C. Carrillo, A. F. Obando Monta\~no, J. Cidr\'as, and E. D\'iaz-Dorado, ``Review of power curve modelling for wind turbines,'' \textit{Renew. Sustain. Energy Rev.}, vol. 21, pp. 572--581, May 2013, doi: 10.1016/j.rser.2013.01.012.

\bibitem{Bastankhah2014}
M. Bastankhah and F. Port\'e-Agel, ``A new analytical model for wind-turbine wakes,'' \textit{Renew. Energy}, vol. 70, pp. 116--123, Oct. 2014, doi: 10.1016/j.renene.2014.01.029.

\end{thebibliography}

\begin{IEEEbiographynophoto}{Prince Solanki}
\TBD{biography: degrees, institutions, research interests}
\end{IEEEbiographynophoto}

\begin{IEEEbiographynophoto}{Pulkit Dwivedi}
\TBD{biography}
\end{IEEEbiographynophoto}

\begin{IEEEbiographynophoto}{Vanita Garg}
\TBD{biography}
\end{IEEEbiographynophoto}

\begin{IEEEbiographynophoto}{Vinay Shukla}
\TBD{biography}
\end{IEEEbiographynophoto}

\end{document}
